\songtitle{El corrido de Carlson}

\tag*{2024}\tag{narcocorrido}\tag{ricochets}\tag{john-carlson}\tag{spanish-lyrics}

\begin{notes}
Based on characters from \emph{Ricochets} by Carlos Thompson.\\
Character John R.~Carlson.\\
Lyrics written by Gemini on a prompt by the author. %
Low interaction.
\end{notes}

\begin{style}
corrido norteño / narcocorrido
\end{style}

\begin{lyrics}

\songpart{Verso 1}
Se las dio de muy astuto,\\
el gringo Carlson mentado,\\
pensó que firmaba un pacto\\
con el diablo de aliado,\\
pero jugó con fuego,\\
y el negocio le ha fallado.

\songpart{Verso 2}
Por un lío de faldas viejas,\\
al cartel quiso usar,\\
les vio la cara de tontos\\
y se puso a celebrar,\\
pero las deudas de sangre\\
siempre te van a cobrar.

\songpart{Coro}
¡Ay, gringo Carlson, qué caro pagaste!\\
La ley de la mafia no sabe perdonar.\\
No te cobraron con oro ni plata,\\
fue tu pequeña la que vio el final.\\
Hoy tu arrogancia se viste de luto,\\
por lo que nunca vas a recuperar.

\songpart{Verso 3}
Dejó la ardiente Miami,\\
donde la dicha perdió,\\
buscando esconder el alma\\
donde no alumbrara el sol.\\
Llegó a la fría Bogotá,\\
con el pecho hecho carbón.

\songpart{Verso 4}
Hoy se le ve en las cantinas,\\
ahogando el frío en alcohol,\\
buscando en cada botella\\
el peso de su error.\\
Pero ni todo el aguardiente\\
le quita ese sinsabor.

\annotation{Tránsito / Hablado}
\emph{(Con música suave y el acordeón llorando)}
> "Y es que en este negocio, gringo...
> la astucia se paga con sangre.
> Ahora gane lo que gane, ya lo perdió todo."

\songpart{Verso Final}
Ahí deambula el gringo Carlson,\\
fantasma entre la neblina,\\
Miami le dio el dinero,\\
Bogotá la cruz divina.

El que la hace la paga,\\
así es la ley de la esquina.

\end{lyrics}

